\section{Relative Trace Formulas, Jacquet--Rallis Matching, and Spectral Analysis}
\label{sec:ch03-rtf-jacquet-rallis-analysis}

To compare automorphic periods on definite unitary groups against periods on
general linear groups---and to control the spectral norms of projected
Poincar\'e series and theta-lift kernels---the proof relies on the
Jacquet--Rallis relative trace formula (RTF) together with a suite of local
harmonic-analytic, complex-geometric, and archimedean representation-theoretic
estimates.
In the Lean formalization, the highest-bottleneck theorems in this part of the
dependency graph form a remarkably self-contained pillar: a complete,
unconditional proof of the Jacquet--Rallis fundamental lemma for all ranks
$r \ge 0$ via Wei Zhang's Fourier-transform method, its descent to the
unitary-group unit lemma and inert principal-congruence matching, H\"ormander's
weighted $L^2$ $\bar{\partial}$-estimates on complex atlases, sharp restriction
and Hermitian curvature bounds, $\mathrm{GL}_3$ mean-square cumulant estimates,
and the spectral geometry of Harish-Chandra $(\mathfrak{g}, K)$-modules.

% -----------------------------------------------------------------------------
\subsection{The Jacquet--Rallis Setting: Orbital Integrals and Invariant Theory}
\label{subsec:ch03-jr-setup}

Let $E/F$ be an unramified quadratic extension of non-archimedean local fields
of characteristic zero with nontrivial involution $x \mapsto \bar{x}$ fixing
$F$ (\leanlink{GlobalTheory.JR.Unram}).
Let $\eta : F \to \{\pm 1\}$ denote the unramified quadratic character
associated to $E/F$, given by $\eta(x) = 1$ if $|x|_F \in |F^{\times}|_F^2$ and
$\eta(x) = -1$ otherwise (\leanlink{GlobalTheory.JR.eta}), and let
$\psi : F \to \mathbb{C}^{\times}$ be a standard additive character of
conductor $\mathcal{O}_F$ (\leanlink{GlobalTheory.JR.IsStdChar}).
On any locally compact Hausdorff space $X$, the space of Schwartz--Bruhat test
functions consists of locally constant, compactly supported functions
$f : X \to \mathbb{C}$ (\leanlink{GlobalTheory.JR.IsTest}).

The infinitesimal Jacquet--Rallis comparison pairs two representations of
reductive groups of rank $r \in \mathbb{N}$:
\begin{enumerate}
  \item \emph{The general linear side} (\leanlink{GlobalTheory.JR.Gl}): the
        ambient space is
        \[
          \mathrm{Gl}_r(k) = \operatorname{Mat}_{r \times r}(k) \times k^r \times (k^r)^{*} \times k,
        \]
        with $g \in \mathrm{GL}_r(k)$ acting by
        $g \cdot (A, v, w, d) = (g A g^{-1},\, g v,\, w g^{-1},\, d)$
        (\leanlink{GlobalTheory.JR.act}).
        A tuple $Y = (A, v, w, d)$ is mapped to its categorical quotient
        invariant (\leanlink{GlobalTheory.JR.inv})
        \[
          \operatorname{inv}(Y)
          = \bigl(c_0(A), \dots, c_{r-1}(A);\, m_0(Y), \dots, m_{r-1}(Y);\, d\bigr)
          \in \mathrm{Inv}_r(k),
        \]
        recording the coefficients of the characteristic polynomial of $A$,
        the Krylov moments $m_i(Y) = w A^i v$ (\leanlink{GlobalTheory.JR.mom}),
        and the scalar $d$.
        Every invariant tuple $a \in \mathrm{Inv}_r(k)$ admits a canonical
        companion-matrix section $\operatorname{sec}(a) \in \mathrm{Gl}_r(k)$
        (\leanlink{GlobalTheory.JR.sec}) whose regularity is detected by the
        Hankel discriminant $\Delta(a) = \det\bigl(m_{i+j}(\operatorname{sec}(a))\bigr)_{0 \le i, j < r}$
        (\leanlink{GlobalTheory.JR.Delta}).
        Bordering $(A, v, w, d)$ into an $(r+1) \times (r+1)$ matrix
        $\widetilde{Y}$ (\leanlink{GlobalTheory.JR.full}) defines the transfer
        factor
        \[
          \omega(Y) = \eta\!\bigl(\det\nolimits_{\mathrm{last}}(Y)\bigr),
        \]
        where $\det_{\mathrm{last}}(Y)$ (\leanlink{GlobalTheory.JR.lastDet}) is
        the Krylov determinant of the row vectors $e_{r+1}^{\top} \widetilde{Y}^i$
        for $0 \le i \le r$.
  \item \emph{The unitary Lie-algebra side} (\leanlink{GlobalTheory.JR.Hn}):
        the ambient space is
        \[
          \mathrm{Hn}_r(F, E) = \operatorname{Herm}_r(E) \times E^r \times F,
        \]
        carrying the action
        $g \cdot (A, v, d) = (g A g^{*},\, g v,\, d)$
        of the unitary group $\mathrm{U}_r(E)$ (\leanlink{GlobalTheory.JR.actU}).
        Embedding $X = (A, v, d) \mapsto (A, v, v^{*}, d) \in \mathrm{Gl}_r(E)$
        defines matching of orbits: $X$ has invariant $a \in \mathrm{Inv}_r(F)$
        (\leanlink{GlobalTheory.JR.HasInv}) precisely when
        $\operatorname{inv}(X^{\mathrm{Gl}}) = a$ in $\mathrm{Inv}_r(E)$.
\end{enumerate}

Endowing $\mathrm{GL}_r(F)$, $\mathrm{U}_r(E)$, $F$, and $E$ with normalized
Haar measures $\mu_G, \mu_U, \mu_F, \mu_E$ assigning volume~$1$ to the
integral maximal compact subgroups $\mathrm{GL}_r(\mathcal{O}_F)$
(\leanlink{GlobalTheory.JR.glO}), $\mathrm{U}_r(\mathcal{O}_E)$
(\leanlink{GlobalTheory.JR.uO}), and the closed unit balls
(\leanlink{GlobalTheory.JR.Meas}), the twisted general linear orbital integral
and the unitary orbital integral are defined by
\[
  O^{\mathrm{GL}}(f', Y)
  = \int_{\mathrm{GL}_r(F)} f'(g \cdot Y)\,\omega(g \cdot Y)\, d\mu_G(g),
  \qquad
  O^{\mathrm{U}}(f, X)
  = \int_{\mathrm{U}_r(E)} f(g \cdot X)\, d\mu_U(g)
\]
(\leanlink{GlobalTheory.JR.orbG}, \leanlink{GlobalTheory.JR.orbU}).
A pair of test functions $(f', f)$ \emph{matches at $a \in \mathrm{Inv}_r(F)$}
(\leanlink{GlobalTheory.JR.MatchAt}) if
$O^{\mathrm{GL}}(f', \operatorname{sec}(a)) = O^{\mathrm{U}}(f, X)$ for every
$X \in \mathrm{Hn}_r(F, E)$ with invariant $a$, and
$O^{\mathrm{GL}}(f', \operatorname{sec}(a)) = 0$ whenever no such $X$ exists;
global matching (\leanlink{GlobalTheory.JR.Match}) requires this identity at
every regular semisimple invariant ($\Delta(a) \neq 0$).
The Jacquet--Rallis fundamental lemma at rank $r$
(\leanlink{GlobalTheory.JR.UnitMatch}) asserts that the characteristic
functions of the integral lattices, $\mathbf{1}_{\mathrm{Gl}_r(\mathcal{O}_F)}$
(\leanlink{GlobalTheory.JR.unitG}) and $\mathbf{1}_{\mathrm{Hn}_r(\mathcal{O}_F, \mathcal{O}_E)}$
(\leanlink{GlobalTheory.JR.unitU}), match.

% -----------------------------------------------------------------------------
\subsection{Fourier Matching and Inductive Proof of the Jacquet--Rallis Fundamental Lemma}
\label{subsec:ch03-jr-proof}

Wei Zhang's strategy proves the Jacquet--Rallis fundamental lemma by showing
first that orbital matching is preserved by Fourier transforms (up to the Weil
constant, which equals~$1$ for the standard self-dual Hermitian space), and
then inducting on the rank $r$ by applying a partial Fourier transform in the
vector variables while exploiting the fact that the unit indicator functions
are self-dual.
In the Lean formalization, this entire architecture is executed without axioms
in \leanmod{Development.AbelianVarietyTheory.Proofs.InertiaEigenvalues}.

The analytic heart of Fourier matching is the rank-one \'etale transfer
theorem: for any \'etale $F$-algebra $K$ with quadratic extension
$L = K \otimes_F E$ (\leanlink{GlobalTheory.JR.Et}), the orbital integrals of
the Weil-representation Fourier transform transfer across matched rank-one
test functions.

\begin{theorem}[Rank-one \'etale orbital Fourier transfer]
\label{thm:ch03-rank-one-transfer}
\leanlink{GlobalTheory.T.etaleOrbitalFourier_rankOne_transfer}
The rank-one \'etale orbital Fourier transfer property
\leanlink{GlobalTheory.JR.RankOne} holds for every unramified quadratic
extension $E/F$ and every reduced finite \'etale $F$-algebra $K$ with
involutive pushout $L = K \otimes_F E$.
\end{theorem}

\begin{proof}
The theorem is established by \leanlink{GlobalTheory.JR.RankOneTransfer.rankOne}.
Using the structural decomposition of \'etale Hermitian pairs
(\leanlink{GlobalTheory.T.etaleHermitianPair_structure_facts}), the Fourier
transform preserves Schwartz--Bruhat functions
(\leanlink{GlobalTheory.T.fourierTransformKProd_locallyConstantCompactSupport}).
To compare the orbital integrals of $\widehat{f}$ and $\widehat{f'}$, one
embeds the rank-one orbital integrals into the Kirillov/Whittaker model of a
Weil representation of $\mathrm{SL}_2(K)$: the truncated Whittaker functional
agrees with the orbital integral
(\leanlink{GlobalTheory.T.truncatedWhittaker_jHyp_eq_orbitalIntegral}),
stabilizes under the Weil representation
(\leanlink{GlobalTheory.T.truncatedWhittaker_weilRep_eventually}), and
satisfies Weyl-element Fourier equivariance
(\leanlink{GlobalTheory.T.jHyp_mem_induced_fourier_weyl},
\leanlink{GlobalTheory.T.weilRep_evalZero_mem_induced}).
Finally, Kirillov vanishing in the limit
(\leanlink{GlobalTheory.T.kirillov_vanishing_of_limit}) forces the difference
between the two transformed orbital integrals to vanish identically.
\end{proof}

Combining rank-one \'etale transfer with the existence of a Weil representation
for any quadratic extension (\leanlink{GlobalTheory.JR.WeilField}) yields
commutation of orbital matching with the full Fourier transforms
\leanlink{GlobalTheory.JR.fourG} and \leanlink{GlobalTheory.JR.fourU} across
every rank $r \in \mathbb{N}$.

\begin{theorem}[All-ranks Fourier matching from a quadratic Weil representation]
\label{thm:ch03-fourier-match}
\leanlink{GlobalTheory.T.jacquetRallis_fourier_matching_of_weil_field}
Assuming the existence of the quadratic Weil representation ($\mathrm{WeilField}$,
\leanlink{GlobalTheory.JR.WeilField}), orbital matching commutes with Fourier
transforms at every rank $r \in \mathbb{N}$: if $(f', f)$ matches on
$\mathrm{Gl}_r(F) \times \mathrm{Hn}_r(F, E)$, then $(\widehat{f'}, \widehat{f})$
also matches (\leanlink{GlobalTheory.JR.FourierMatch}).
\end{theorem}

\begin{proof}
The statement reduces immediately to
\leanlink{GlobalTheory.JR.fourierMatch_of_weilField}, which feeds four
ingredients into the general Weil-representation transfer theorem
\leanlink{GlobalTheory.T.jacquetRallisFourierMatch_of_weilRepresentation}:
first, lifting the field-level Weil representation to finite \'etale algebras
(\leanlink{GlobalTheory.T.weilRepresentationEtale_of_quadraticField}); second,
Theorem~\ref{thm:ch03-rank-one-transfer}
(\leanlink{GlobalTheory.T.etaleOrbitalFourier_rankOne_transfer}); and third,
the evaluation of the Weil constant of the standard Hermitian space as $1$
(\leanlink{GlobalTheory.T.weilConstantStandardHermitian_eq_one}).
\end{proof}

We now arrive at the centerpiece theorem of the local relative trace formula:
the unconditional Jacquet--Rallis fundamental lemma in all ranks.

\begin{theorem}[Jacquet--Rallis fundamental lemma across all ranks]
\label{thm:ch03-jr-unit-match-all}
\leanlink{GlobalTheory.T.jacquetRallisUnitMatch_all_ranks}
For every rank $r \in \mathbb{N}$, every unramified quadratic extension $E/F$
of characteristic-zero non-archimedean local fields, and every normalized
measure datum \leanlink{GlobalTheory.JR.Meas}, the unit indicator functions
$\mathbf{1}_{\mathrm{Gl}_r(\mathcal{O}_F)}$ and
$\mathbf{1}_{\mathrm{Hn}_r(\mathcal{O}_F, \mathcal{O}_E)}$ have matching
orbital integrals at all regular semisimple invariants
(\leanlink{GlobalTheory.JR.UnitMatch}).
\end{theorem}

\begin{proof}
By \leanlink{GlobalTheory.JR.unitMatch_all}, the theorem is deduced from the
conditional reduction
\leanlink{GlobalTheory.T.jacquetRallisUnitMatch_of_weilRep_quadratic} applied
to the explicit construction of the Weil representation for quadratic local
fields (\leanlink{GlobalTheory.T.weilRepresentation_quadraticField_exists},
proved via \leanlink{GlobalTheory.JR.WeilQuad.weilField_main}).
Inside \leanlink{GlobalTheory.JR.unitMatch_of_weilField}, one invokes
Theorem~\ref{thm:ch03-fourier-match}
(\leanlink{GlobalTheory.T.jacquetRallis_fourier_matching_of_weil_field}) and
proceeds by induction on $r$:
\begin{itemize}
  \item \emph{Base case ($r = 0$)}: at rank zero, the matching identity
        \leanlink{GlobalTheory.JR.SizeOne.unitMatch_zero} is a direct
        evaluation of unit-ball volumes.
  \item \emph{Inductive step ($r \to r + 1$)}: assuming unit matching at rank
        $r$ and Fourier matching at rank $r + 1$, the inductive step
        \leanlink{GlobalTheory.JR.UnitStep.unitMatch_succ} verifies matching on
        the dense locus of invariants with separable characteristic polynomial
        (\leanlink{GlobalTheory.T.jacquetRallis_match_of_dense_separable}).
        When some eigenvalue or invariant has norm at least~$1$, orbital
        descent reduces matching to rank $r$ via
        \leanlink{GlobalTheory.JR.Step.unitMatchAt_of_one_le_norm}; on the
        complementary small-norm locus, applying the partial Fourier transform
        (\leanlink{GlobalTheory.T.partialFourier_preserves_test_fixes_unit}),
        which fixes the unit test functions, transfers the vanishing of the
        difference of orbital integrals from the large-norm locus to the
        entire regular semisimple locus.
\end{itemize}
\end{proof}

In the global proof, Theorem~\ref{thm:ch03-jr-unit-match-all} is immediately
lifted from the Lie algebra to unitary groups via Cayley transforms and
involutive algebras, yielding two fundamental applications that enter the
base-change and trace-formula comparisons.

\begin{corollary}[Unitary-group unit lemma and inert principal-congruence matching]
\label{cor:ch03-unit-lemma-cong}
\leanlink{GlobalTheory.UnitLemma.unitLemma_identity_unconditional}
\begin{enumerate}
  \item For every unramified quadratic local datum $K$ and every pair of
        involutive algebras $B_1, B_2 \in \operatorname{InvAlg}(K)$ with
        unitary elements $g_i \sigma(g_i) = 1$ whose characteristic polynomials
        have separable product $P_1 P_2$, the unitary-group unit-lemma identity
        holds unconditionally
        (\leanlink{GlobalTheory.UnitLemma.unitLemma_identity_unconditional}).
  \item Let $E$ be an unramified CM number field, $w$ a finite place of $E$
        that is not split over its totally real subfield, $n = a + b$, and
        $k \ge 1$.
        Then the indicator function of the $k$-th principal congruence subgroup
        of $\mathrm{U}_n(E_w)$ matches the rescaled product norm-congruence
        test function on $\mathrm{GL}_a(E_w) \times \mathrm{GL}_b(E_w)$ with
        explicit transfer scalar $(-\sqrt{q_w})^{-k a b}$
        (\leanlink{GlobalTheory.CongMatch.match_principal_congruence_inert}).
\end{enumerate}
\end{corollary}

\begin{proof}
Both parts factor directly through
Theorem~\ref{thm:ch03-jr-unit-match-all}
(\leanlink{GlobalTheory.T.jacquetRallisUnitMatch_all_ranks}): the Lie-algebra
matching is transferred to the unitary-group Lie algebra and then to unitary
groups at every residue characteristic via
\leanlink{GlobalTheory.T.unitLemma_lieAlgebra_unitaryGroup_of_jacquetRallis} and
\leanlink{GlobalTheory.T.unitLemma_unitary_group_every_residue_of_lie}.
\end{proof}

% -----------------------------------------------------------------------------
\subsection{Weighted $L^2$ $\bar{\partial}$-Estimates and Chart Restriction on Reductive Models}
\label{subsec:ch03-dbar-charts}

Alongside the non-archimedean orbital-integral matching, the proof constructs
holomorphic sections of Hermitian line bundles with quantitative $L^2$ control
at the archimedean places (essential for arithmetic Hilbert--Samuel, Moriwaki
bigness, and archimedean period bounds) and glues line-bundle isomorphisms on
integral models of reductive and abelian schemes.

The complex-analytic core is H\"ormander's weighted $L^2$ $\bar{\partial}$-theorem
on complex manifolds covered by a weighted atlas $A$
(\leanmod{ReciprocityTheory.L2DbarEstimates.PlurisubharmonicCharts}).
Given a plurisubharmonic potential $\Phi$ whose Levi form is nonnegative on
every local chart, and a smooth section $F_0$ whose frame quotient $u_0$
satisfies the pointwise Levi-form domination inequality
\[
  \bigl\lvert \bar{\partial} u_0(x) \cdot \eta \bigr\rvert^2
  \;\le\;
  h_{\mathrm{eff}}(x)\,\operatorname{Re}\bigl(\operatorname{Levi}(\Phi)(x; \eta)\bigr),
\]
solving $\bar{\partial} v = \bar{\partial} u_0$ in $L^2(e^{-\Phi})$ produces a
holomorphic correction with controlled weighted $L^2$ norm.

\begin{theorem}[Weighted $L^2$ $\bar{\partial}$-estimate on complex atlases]
\label{thm:ch03-l2-dbar}
\leanlink{ReciprocityTheory.L2DbarEstimates.Y.hy}
Let $n \in \mathbb{N}$ and let $Y$ be a second-countable Hausdorff Borel space.
Given the flat-space H\"ormander estimate and null-extension inputs
$(h_H, h_E, h_N)$, the global weighted $L^2$ $\bar{\partial}$-estimate holds on
every weighted complex atlas on $Y$: whenever the weighted integral
$\int_Y h(y)\, e^{-\psi(y)}\, d\mu(y)$ is finite and the modified weight $\Phi$
has nonnegative Levi form dominating $|\bar{\partial} u_0|^2$, there exists a
holomorphic solution whose weighted $L^2$ norm is bounded by that integral.
\end{theorem}

\begin{proof}
In the formalization, one passes to the open Stein complement $M$ where the
local denominator section $w_0$ is nonvanishing and transports the problem to
a flat complex chart via local Levi-form flattening
(\leanlink{ReciprocityTheory.L2DbarEstimates.Y.levi_flat_loc}).
The Levi-form identity, integral transport
(\leanlink{ReciprocityTheory.L2DbarEstimates.Y.lintegral_flat}), and
holomorphicity equivalence
(\leanlink{ReciprocityTheory.L2DbarEstimates.Y.holo_flat_iff}) convert the
atlas hypotheses into the flat-space H\"ormander estimate $h_H$; the resulting
holomorphic solution on $M$ extends across the analytic hypersurface by the
extension inputs $h_E, h_N$.
\end{proof}

On the arithmetic-geometric side of \leanmod{ReciprocityTheory.ReductiveGroups},
line-bundle descent across discrete-valuation charts and toroidal
compactifications ensures that integral structures on automorphic bundles are
compatible across overlaps and group isomorphisms.
Specifically, for line bundles $A, B$ on a scheme $Y$ with a generating
section $u$ of $B \otimes A^{-1}$, the candidate chart-restricted sections
\leanlink{ReciprocityTheory.ChartRestriction.G.DVR.isoCand} agree on all
triple overlaps (\leanlink{ReciprocityTheory.ChartRestriction.G.DVR.isoCand_compat})
and glue to a global section \leanlink{ReciprocityTheory.ChartRestriction.G.DVR.secOfIso}.
Similarly, in the toroidal compactification layer, the group-theoretic
toroidal property is invariant under monoidal scheme isomorphisms
(\leanlink{ReciprocityTheory.ChartRestriction.G.Geo.Tor.isGt_of_iso}).

% -----------------------------------------------------------------------------
\subsection{Sharp Restriction, Hermitian Curvature, and $\mathrm{GL}_3$ Mean-Square Estimates}
\label{subsec:ch03-sharp-curvature-gl3}

Once orbital matching and global section existence are in place, controlling
the spectral side of the relative trace formula and the norms of projected
Poincar\'e series requires three quantitative analytic packages: sharp
restriction bounds on cuspidal spans, Hermitian curvature and Fourier-decay
evaluations, and $\mathrm{GL}_3$ mean-square / cumulant estimates.

\paragraph{Sharp restriction and cusp-span norm bounds.}
In \leanmod{Formalism.SharpRestriction}, the spectral parameters of the
automorphic family are encoded by the word-indexed structure
\leanlink{Formalism.MainNotions.S.Draft.SData}, which spaces the Harish-Chandra
parameters uniformly by a positive cloud gap $4\Delta$
(\texttt{cloud\_gap}) and encodes each word $w$ as a binary pattern of
half-integer shifts by $\Delta$ (\texttt{u\_letter}).
Each word $u$ carries a finite-dimensional cuspidal space $C_u$ embedded into a
common Hilbert space $H$ by a continuous projection $P_u : C_u \to H$.
To compare inner products against local Bessel integrals,
\leanmod{Formalism.SharpRestriction.CuspSpanNormBounds.RelevantZeroGramTerms}
formalizes the formal degree $d(\sigma) > 0$ of an irreducible square-integrable
unitary representation $\sigma$ of an indefinite unitary group $U(J_\varepsilon)$
(\leanlink{Formalism.MainNotions.S.IsFdeg}, \leanlink{Formalism.MainNotions.S.fdegOf})
via the Schur orthogonality identity
\[
  \int_{U(J_\varepsilon)} \langle w, \sigma(g) v \rangle\,
  \overline{\langle w', \sigma(g) v' \rangle}\, d\mu(g)
  \;=\;
  \frac{\langle v', v \rangle\, \langle w, w' \rangle}{d(\sigma)}.
\]
At the finite places, the integral moduli charts
\leanlink{Formalism.MainNotions.ModuliIntegralCharts.ExtA} and the totient
Galois-class host derivation \leanlink{Formalism.HC.D7.hostDer} (combining
inertia and discriminant control on the host CM field) interface with the
$p$-adic logarithmic theta series in
\leanmod{Formalism.SharpRestriction.GlobalSchurBounds.GL3Theta}.

\paragraph{Hermitian curvature and explicit first-kernel Fourier decay.}
In \leanmod{Formalism.HermitianCurvature}, archimedean Fourier decay is
formulated using an abstract differential algebra
\leanlink{Formalism.FourierDecayAlt.MA.DAlg} equipped with commuting Euler
derivations satisfying the Leibniz rule on coordinates $x_v$ and inverting the
curvature determinant $g_T = z \bar{z} - T$.
On the open smooth domain $Q_{\mathrm{set}}(r, T)$
(\leanlink{Formalism.FourierDecayAlt.MA.Qset_open}), both the coordinate
functions (\leanlink{Formalism.FourierDecayAlt.MA.coordE_smooth}) and their
Wirtinger derivatives (\leanlink{Formalism.FourierDecayAlt.MA.wirt_smooth}) are
smooth, allowing iterated integration by parts against the oscillatory phase.
At non-archimedean places, the explicit first-kernel calculation evaluates the
inner exponential sum attached to a local datum $(K, R, \psi)$ and integral
matrices $Z_1 \in \operatorname{Mat}_m(K)$, $Z_2 \in \operatorname{Mat}_{m'}(K)$.

\begin{proposition}[Explicit first-kernel inner exponential sum]
\label{prop:ch03-first-kernel-sum}
\leanlink{Formalism.FourierDecayAlt.P.A5x8.inner_sum}
Let $R \subset K$ be a local subring and $\psi : K \to \mathbb{C}^{\times}$ an
additive character satisfying the local datum axioms.
For $N = m + 1$, $N' = m' + 1$, integral matrices $Z_1, Z_2$ with unit
determinants, and a unit $c \in K^{\times}$, the first-kernel inner exponential
sum evaluates explicitly via the block-matrix factorization and injective
indexing map $\iota$.
\end{proposition}

\begin{proof}
The proof reduces the exponential sum over the bordered matrix domain to a
supported sum along the injective map
\leanlink{Formalism.FourierDecayAlt.P.A5x8.iota_injective}, using the support
localization \leanlink{Formalism.FourierDecayAlt.P.A5x8.tTerm_supp} and the
exact term evaluation \leanlink{Formalism.FourierDecayAlt.P.A5x8.tTerm_iota}.
\end{proof}

\paragraph{$\mathrm{GL}_3$ mean-square and Hermitian-lattice cumulant bounds.}
In \leanmod{Formalism.GL3MeanSquares}, the off-diagonal mean-square bounds for
$\mathrm{GL}_3$ theta lifts depend on two combinatorial-arithmetic ingredients:
first, an exact inversion-count identity for hook-formula indices
(\leanlink{Formalism.HermitianLattices.CumulantEstimateDefs.CumulantEstimates.inner})
controlling the Hilbert--Schmidt half-bound in the cumulant expansion; second,
the uniqueness and Frobenius-orbit stability of conjugate ideal factorizations
$D = B \cdot \bar{B}$ in the ring of integers of the host CM field
(\leanlink{Formalism.HermitianLattices.FiltrationBoundaryCriterion.PfN.principalStable}),
paired with involutive symmetry of uniformizer divisibility for unimodular
Hermitian genus representatives
(\leanlink{Formalism.HermitianLattices.Ext.L3.U0.Dv.star}).

% -----------------------------------------------------------------------------
\subsection{Spectral Geometry and Harish-Chandra Modules for Unitary Groups}
\label{subsec:ch03-spectral-harish-chandra}

The final analytic pillar, formalized in
\leanmod{CompanionTheory.HarishChandraModules} and
\leanmod{CompanionTheory.SpectralGeometry}, supplies the archimedean
representation theory of indefinite unitary groups $U(p, q)$ needed to isolate
discrete-series packets, transfer differential-operator identities between
smooth and $K$-finite vectors, and control low $K$-type multiplicities.

\begin{theorem}[Iwasawa decomposition and admissible $(\mathfrak{g}, K)$-module transfer]
\label{thm:ch03-hc-transfer}
\leanlink{CompanionTheory.D.gk.admitted.transfer_holds}
For indefinite unitary groups $U(J_\varepsilon)$ of every signature
$\varepsilon : \{1, \dots, n\} \to \{\pm 1\}$:
\begin{enumerate}
  \item \emph{Iwasawa decomposition}
        (\leanlink{CompanionTheory.D.gk.admitted.iwasawa_decomposition_unitary_group}):
        every unitary group $U(J_\varepsilon)$ admits a global Iwasawa
        decomposition constructed via signed Cholesky factorization.
  \item \emph{Admissible smooth-to-$K$-finite transfer}
        (\leanlink{CompanionTheory.D.gk.admitted.transfer_holds}):
        standard polynomial differential-operator relations on $K$-finite
        vectors transfer cleanly to admissible smooth vectors via uniform
        approximation and norm bounds.
  \item \emph{Minimal-eigenspace integrality}
        (\leanlink{CompanionTheory.D.gk.admitted.goodPI_of}):
        from small-parameter integrality, common eigenvectors exist and
        generate the full minimal $K$-type eigenspace.
\end{enumerate}
\end{theorem}

\begin{proof}
In the Lean formalization, the transfer theorem
\leanlink{CompanionTheory.D.gk.admitted.transfer_holds} combines the
approximation lemma
\leanlink{CompanionTheory.D.gk.admitted.PrismTransfer.exists_approx}, the
operator-norm bound
\leanlink{CompanionTheory.D.gk.admitted.PrismTransfer.exists_bound}, and
compatibility with standard polynomials
\leanlink{CompanionTheory.D.gk.admitted.map_stdPoly}.
\end{proof}

To bound the multiplicity of minimal $K$-types in discrete-series
representations, \leanmod{CompanionTheory.HarishChandraModules.LowKTypeMultiplicity}
realizes the generalized Verma module algebraically as a polynomial module and
analyzes its Koszul creation-operator complex.

\begin{theorem}[Generalized Verma module realization and Koszul creation-operator acyclicity]
\label{thm:ch03-verma-koszul}
\leanlink{CompanionTheory.D.LK.generalized_verma_module_polynomial}
For all ranks $n, m \in \mathbb{N}$:
\begin{enumerate}
  \item The generalized Verma module of a discrete-series Harish-Chandra module
        admits an explicit polynomial-ring realization in which the Lie-algebra
        action acts by first-order polynomial differential operators
        (\leanlink{CompanionTheory.D.LK.generalized_verma_module_polynomial}).
  \item The associated graded module is free for every $(n, m)$
        (\leanlink{CompanionTheory.D.LK.free_all}, proved by
        \leanlink{CompanionTheory.D.LK.free_module_induction}).
  \item The Koszul creation-operator differential satisfies $d^2 = 0$, has
        bounded support
        (\leanlink{CompanionTheory.D.LK.differential_bounded_support}), is
        acyclic (\leanlink{CompanionTheory.D.LK.acycAbs_of}), and generates the
        creation-operator module unconditionally
        (\leanlink{CompanionTheory.D.LK.hGenCO}, proved via
        \leanlink{CompanionTheory.D.LK.genCO_of}).
\end{enumerate}
\end{theorem}

\begin{proof}
The polynomial realization
\leanlink{CompanionTheory.D.LK.generalized_verma_module_polynomial} is verified
by checking the abstract $\mathfrak{gl}$-commutation relations
(\leanlink{GlobalTheory.CompanionTheory.DSVerma.glRel_abstract}) together with
the differential commutativity identities
\leanlink{GlobalTheory.CompanionTheory.DSVerma.dd_comm} and
\leanlink{GlobalTheory.CompanionTheory.DSVerma.dd_yy}.
\end{proof}

Finally, \leanmod{CompanionTheory.SpectralGeometry} bridges these archimedean
$(\mathfrak{g}, K)$-module properties back to the spectral geometry of
automorphic bundles and stable lattices: the low-$K$-type span property holds
unconditionally (\leanlink{GlobalTheory.CompanionTheory.KLowSpan.lowSpan_holds})
and produces the canonical highest-weight map
(\leanlink{GlobalTheory.CompanionTheory.DSICHWMap.hwMap_of_lowSpan}), while
flat Hilbert~90 yields rigidified descent of line bundles on arithmetic models
(\leanlink{GlobalTheory.CompanionTheory.RigDesc.rigDescent}).
Together with the all-ranks Jacquet--Rallis fundamental lemma
(Theorem~\ref{thm:ch03-jr-unit-match-all}), the weighted $L^2$
$\bar{\partial}$-estimate (Theorem~\ref{thm:ch03-l2-dbar}), and the
analytic restriction and curvature bounds of
Subsection~\ref{subsec:ch03-sharp-curvature-gl3}, this completes the
automorphic-analytic machinery supporting the global proof of the Generalized
Riemann Hypothesis.
